% Web source: docs/05statistical_mechanics/01/01_040.md
\addcontentsline{toc}{section}{1-4　ボルツマンエントロピー}

\begin{mondai}{1-4}{ボルツマンエントロピー}{標準}
孤立した容器の中に，互いに相互作用をしない区別可能な $N$ 個の粒子からなる気体が存在している。容器の内部は中央の仮想的な境界によって左側と右側の2つの等しい領域に分かれており，粒子はどちらの領域にも等しい確率で入ることができるものとする。左側の領域に存在する粒子数を $n$，右側の領域に存在する粒子数を $N-n$ とするとき，以下の問いに答えよ。

\begin{enumerate}
  \item 左側の領域にちょうど $n$ 個の粒子が存在するときの微視的状態数 $\Omega(n)$ を求めよ。

  \item ボルツマンのエントロピー公式 $S(n) = k_B \log \Omega(n)$ にスターリングの公式 $\log X! \simeq X \log X - X$ を適用し，エントロピー $S(n)$ を $N$ と $n$ を用いて表せ。

  \item エントロピー $S(n)$ を $n$ について微分し，最大値をとる $n$ を求めよ。さらに，$n=N/2+\delta n$ と置いて $S(n)$ を $\delta n$ の2次まで展開し，状態数の比が
  \begin{equation*}
  \frac{\Omega(N/2+\delta n)}{\Omega(N/2)} \simeq \exp\left[-\frac{2(\delta n)^2}{N}\right]
  \end{equation*}
  となることを示せ。
\end{enumerate}
\end{mondai}

\solutionheading
\begin{solutionparts}

\solutionpart{(1)}
$N$ 個の区別可能な粒子の中から，左側の領域に入る $n$ 個の粒子を選ぶ組合せの数が，指定された配置の微視的状態数 $\Omega(n)$ に対応します。したがって，求める状態数は二項係数を用いて以下のように表されます。
\begin{equation*}
\Omega(n) = \frac{N!}{n!(N-n)!}
\end{equation*}

\solutionpart{(2)}
ボルツマンのエントロピー公式に (1) の結果を代入し，対数の性質を用いて展開します。
\begin{equation*}
S(n) = k_B \log \left[ \frac{N!}{n!(N-n)!} \right] = k_B \left[ \log N! - \log n! - \log (N-n)! \right]
\end{equation*}
ここで，各項にスターリングの公式を適用します。
\begin{equation*}
S(n) \simeq k_B \left[ (N \log N - N) - (n \log n - n) - \{ (N-n) \log (N-n) - (N-n) \} \right]
\end{equation*}
括弧内の定数項 $-N + n + (N-n) = 0$ となって相殺されるため，エントロピーは以下のようになります。
\begin{equation*}
S(n) \simeq k_B \left[ N \log N - n \log n - (N-n) \log (N-n) \right]
\end{equation*}

\solutionpart{(3)}
エントロピー $S(n)$ が最大となる条件を調べるため，$n$ に関する1階微分を計算して 0 とおきます。
\begin{equation*}
\frac{dS(n)}{dn} = k_B \left[ - \left( \log n + n \cdot \frac{1}{n} \right) - \left( - \log (N-n) + (N-n) \cdot \frac{1}{N-n} \cdot (-1) \right) \right]
\end{equation*}
\begin{equation*}
\frac{dS(n)}{dn} = k_B \left[ - \log n - 1 + \log (N-n) + 1 \right] = k_B \log \left( \frac{N-n}{n} \right)
\end{equation*}
極値をとる条件 $\dfrac{dS(n)}{dn} = 0$ より，以下の関係式が得られます。
\begin{equation*}
\frac{N-n}{n} = 1
\end{equation*}
\begin{equation*}
n = \frac{N}{2}
\end{equation*}
2階微分を計算すると $\dfrac{d^2S}{dn^2} = -k_B \left( \dfrac{1}{n} + \dfrac{1}{N-n} \right) < 0$ となるため，$n = N/2$ のときにエントロピーは確かに最大値をとります。
\par
$n=N/2$ におけるエントロピーと2階微分は
\begin{equation*}
S\left(\frac{N}{2}\right)=Nk_B\log 2,\qquad
\left.\frac{d^2S}{dn^2}\right|_{n=N/2}=-\frac{4k_B}{N}
\end{equation*}
です。したがって，$n=N/2+\delta n$ として2次まで展開すると，
\begin{equation*}
S\left(\frac{N}{2}+\delta n\right)
\simeq Nk_B\log 2-\frac{2k_B}{N}(\delta n)^2
\end{equation*}
となります。$S=k_B\log\Omega$ を用いれば，状態数の比は
\begin{equation*}
\frac{\Omega(N/2+\delta n)}{\Omega(N/2)}
=\exp\left[\frac{S(N/2+\delta n)-S(N/2)}{k_B}\right]
\simeq \exp\left[-\frac{2(\delta n)^2}{N}\right]
\end{equation*}
です。したがって，粒子数のずれの典型的な大きさは $\sqrt{N}$ のオーダーであり，相対的なずれは $1/\sqrt{N}$ のオーダーで0に近づきます。

\end{solutionparts}
